\documentclass[11pt]{article}

\usepackage[margin=1in]{geometry}
\usepackage{amsmath,amssymb,amsthm}
\usepackage{graphicx}
\usepackage{algorithm}
\usepackage{algpseudocode}
\usepackage{enumitem}
\usepackage{hyperref}

\DeclareMathOperator{\argmax}{arg\,max}
\newtheorem{definition}{Definition}
\newtheorem{lemma}{Lemma}
\newtheorem{proposition}{Proposition}
\newtheorem{theorem}{Theorem}
\newtheorem{corollary}{Corollary}
\newtheorem{remark}{Remark}

\title{Validity of Anti-Monotone Bound-and-Prune Search\\
for the Maximally Violating Conjunctive Subgroup\\
\large (unifying MSD, SPSF and FPSF)}
\author{}
\date{}

\begin{document}
\maketitle

\begin{abstract}
We formalize and prove the correctness of a level-wise (breadth-first, increasing-order)
branch-and-bound search over conjunctive subgroups that finds the subgroup of maximal
violation for the three intersectional-bias measures Maximum Subgroup Discrepancy (MSD)
\cite{msd}, Statistical Parity Subgroup Fairness (SPSF) and False Positive Subgroup
Fairness (FPSF) \cite{kearns18}. The search prunes a subgroup, and \emph{every} subgroup
obtained from it by specifying further attributes, as soon as a cheaply computable
upper bound on its violation falls below the best violation already found. We isolate
the single abstract property (an \emph{anti-monotone optimistic bound}) that makes this
sound for all three measures, instantiate it exactly as the two bound formulas quoted
in the problem statement, show that a canonical (order-restricted) generation scheme
produces every subgroup exactly once — so the same pruning check automatically covers
every path by which a ``pointless'' subgroup could otherwise be reconstructed from a
different parent — and prove that the returned value and subgroup are always globally
optimal.
\end{abstract}

\section{Setting: the conjunctive subgroup lattice}

Let $A=\{a_1,\dots,a_d\}$ be a finite set of discrete protected attributes, attribute
$a$ taking values in a finite domain $\mathrm{Val}(a)$. An \emph{attribute-value pair}
(a \emph{literal}) is a statement $a\!=\!v$ with $a \in A$, $v\in\mathrm{Val}(a)$. A
\emph{subgroup} is a set of literals using every attribute at most once,
\[
S=\{a_{i_1}\!=\!v_1,\dots,a_{i_k}\!=\!v_k\},\qquad a_{i_1},\dots,a_{i_k}\in A \text{ pairwise distinct.}
\]
Write $\mathrm{attr}(S)\subseteq A$ for the attributes it specifies and $|S|=k$ for its
\emph{order}. $S_A$ denotes the set of \emph{all} such subgroups over all subsets of
$A$ — exactly the family of ``intersections of attribute-value pairs for all subsets
of $A$'' used to define $\mathcal{S}_A$ for MSD, SPSF and FPSF alike. For a point $x$,
$S(x):=\prod_{(a=v)\in S}\mathbf{1}[x_a=v]$, and $X_S:=\{x : S(x)=1\}$.

For $S,S'\in S_A$, write $S \preceq S'$ (``$S'$ \emph{refines} $S$'', or $S$ is a
\emph{sub-conjunction}/\emph{ancestor} of $S'$) iff $S\subseteq S'$ as sets of literals.
$(S_A,\preceq)$ is a partial order with least element $\emptyset$ (the whole
population), and it is exactly the order in which the rough algorithm walks the search
space: $|S|=1$ before $|S|=2$ before $|S|=3,\dots$, i.e.\ any $S\preceq S'$ with
$S\neq S'$ satisfies $|S|<|S'|$.

\begin{lemma}[Coverage monotonicity]\label{lem:cov}
For any (possibly sub-probability) measure $\mu$ on the input space and any
$S\preceq S'$ in $S_A$, $X_{S'}\subseteq X_S$ and hence $\mu(X_{S'})\le \mu(X_S)$.
\end{lemma}
\begin{proof}
$S\subseteq S'$ means every constraint defining $X_S$ is also imposed by $S'$, so
$X_{S'}=\bigcap_{\ell\in S'}X_\ell \subseteq \bigcap_{\ell\in S}X_\ell = X_S$.
Monotonicity of measures under inclusion gives $\mu(X_{S'})\le\mu(X_S)$.
\end{proof}

This is the only structural fact used below: \emph{specifying more attribute-value
pairs can only shrink the covered population, under every distribution at once.}
Applying it with $\mu=P(\cdot,y=0)$ (a sub-probability measure) shows it also
governs FPSF's restriction to the negative class.

\section{The three violation measures}

\begin{definition}[MSD \cite{msd}]
Given two distributions $D_1,D_2$ on $X$,
\[
\mathrm{MSD}(D_1,D_2;A)=\sup_{S\in S_A}\mathrm{SD}(S),\qquad
\mathrm{SD}(S):=\bigl|P_{D_1}(X_S)-P_{D_2}(X_S)\bigr|.
\]
\end{definition}

\begin{definition}[SPSF \cite{kearns18}]
Given a distribution $D$ and classifier $h$, let $p_{\mathrm{base}}:=P(h(x)=1)$. Then
\[
\mathrm{SPSF}(h,S):=P(X_S)\cdot\bigl|p_{\mathrm{base}}-P(h(x)=1\mid X_S)\bigr|.
\]
$h$ satisfies $\gamma$-SPSF w.r.t.\ $S_A$ iff $\mathrm{SPSF}(h,S)\le\gamma$ for every
$S\in S_A$; detection asks for $S^\star\in\argmax_{S\in S_A}\mathrm{SPSF}(h,S)$.
\end{definition}

\begin{definition}[FPSF \cite{kearns18}]
With $p_{\mathrm{base},0}:=P(h(x)=1\mid y=0)$,
\[
\mathrm{FPSF}(h,S):=P(X_S,y=0)\cdot\bigl|p_{\mathrm{base},0}-P(h(x)=1\mid X_S,y=0)\bigr|.
\]
\end{definition}

The unifying paper ties the three together: $\mathrm{FPSF}(h,S)=P(y=0)\cdot
\mathrm{SPSF}_{D\mid y=0}(h,S)$ (FPSF \emph{is} SPSF evaluated on the sub-population
$\{y=0\}$), and, writing $D_{h=1},D_{h=0}$ for the label-conditional distributions
induced by $h$,
\[
\mathrm{SD}\bigl(S;D_{h=1},D_{h=0}\bigr)=\frac{\mathrm{SPSF}(h,S)}{P(h(x)=1)\,P(h(x)=0)},
\]
a proportionality with a constant \emph{independent of $S$}, hence
\[
\argmax_{S\in S_A}\mathrm{SPSF}(h,S)=\argmax_{S\in S_A}\mathrm{SD}(S;D_{h=1},D_{h=0})
\]
(Corollary~1.1, ``Equivalence of SPSF and MSD''). This is why one abstract algorithm,
instantiated three ways below, suffices for all three measures, and why finding the
arg-max for one is exactly as good as finding it for the others.

\section{An anti-monotone optimistic bound}

\begin{definition}[Valid anti-monotone bound]
Let $V:S_A\to\mathbb{R}_{\ge0}$ be a violation function. We call
$U:S_A\to\mathbb{R}_{\ge0}$ a \emph{valid anti-monotone bound} for $V$ if
\begin{enumerate}[label=(B\arabic*)]
\item \textbf{Validity:} $V(S)\le U(S)$ for every $S\in S_A$;
\item \textbf{Anti-monotonicity:} $S\preceq S' \implies U(S')\le U(S)$.
\end{enumerate}
\end{definition}

\begin{lemma}[Bound transfer]\label{lem:transfer}
If $U$ is a valid anti-monotone bound for $V$, then for every $S\preceq S'$ in $S_A$,
$V(S')\le U(S)$.
\end{lemma}
\begin{proof}
$V(S')\overset{(B1)}{\le}U(S')\overset{(B2)}{\le}U(S)$.
\end{proof}

Lemma~\ref{lem:transfer} is the whole point: a bound computed \emph{at a small
subgroup $S$}, using only $S$'s own (cheap) coverage, dominates the exact violation of
\emph{every subgroup obtainable from $S$ by adding more attribute-value pairs} —
regardless of which pairs, how many, or in what order they are added. We now build such
a $U$ for each measure; the two resulting formulas are exactly the two bullet points of
the informal description.

\begin{proposition}[Bound for SPSF]\label{prop:spsf}
Let $\bar\delta_{\mathrm{SP}}:=\max(p_{\mathrm{base}},1-p_{\mathrm{base}})$ and
$U_{\mathrm{SP}}(S):=P(X_S)\cdot\bar\delta_{\mathrm{SP}}$. Then $U_{\mathrm{SP}}$ is a
valid anti-monotone bound for $\mathrm{SPSF}(h,\cdot)$.
\end{proposition}
\begin{proof}
\emph{(B1)} Since $P(h=1\mid X_S)\in[0,1]$, the farthest this conditional probability
can be from the fixed constant $p_{\mathrm{base}}$ while staying in $[0,1]$ is
$\max(p_{\mathrm{base}},1-p_{\mathrm{base}})=\bar\delta_{\mathrm{SP}}$, so
$|p_{\mathrm{base}}-P(h=1\mid X_S)|\le\bar\delta_{\mathrm{SP}}$; multiplying by
$P(X_S)\ge0$ gives $\mathrm{SPSF}(h,S)\le U_{\mathrm{SP}}(S)$.
\emph{(B2)} $\bar\delta_{\mathrm{SP}}$ is a global constant (it does not depend on
$S$); by Lemma~\ref{lem:cov}, $S\preceq S'\Rightarrow P(X_{S'})\le P(X_S)$, hence
$U_{\mathrm{SP}}(S')\le U_{\mathrm{SP}}(S)$.
\end{proof}

\begin{proposition}[Bound for FPSF]\label{prop:fpsf}
Let $\bar\delta_{\mathrm{FP}}:=\max(p_{\mathrm{base},0},1-p_{\mathrm{base},0})$ and
$U_{\mathrm{FP}}(S):=P(X_S,y=0)\cdot\bar\delta_{\mathrm{FP}}$. Then $U_{\mathrm{FP}}$
is a valid anti-monotone bound for $\mathrm{FPSF}(h,\cdot)$.
\end{proposition}
\begin{proof}
Identical to Proposition~\ref{prop:spsf}, applying Lemma~\ref{lem:cov} to the
sub-probability measure $\mu(\cdot)=P(\cdot,y=0)$ in place of $P(\cdot)$.
\end{proof}

\begin{proposition}[Bound for MSD]\label{prop:msd}
Let $U_{\mathrm{MSD}}(S):=\max\bigl(P_{D_1}(X_S),P_{D_2}(X_S)\bigr)$. Then
$U_{\mathrm{MSD}}$ is a valid anti-monotone bound for $\mathrm{SD}(\cdot)$.
\end{proposition}
\begin{proof}
\emph{(B1)} For any $a,b\ge0$, $a-b\le a\le\max(a,b)$ and $b-a\le b\le\max(a,b)$, so
$|a-b|\le\max(a,b)$. Taking $a=P_{D_1}(X_S),\,b=P_{D_2}(X_S)$ gives
$\mathrm{SD}(S)=|a-b|\le\max(a,b)=U_{\mathrm{MSD}}(S)$.
\emph{(B2)} By Lemma~\ref{lem:cov} applied separately to $D_1$ and $D_2$,
$S\preceq S'$ implies $P_{D_1}(X_{S'})\le P_{D_1}(X_S)$ and
$P_{D_2}(X_{S'})\le P_{D_2}(X_S)$; the max of two coordinatewise-dominated
nonnegative numbers is itself dominated, so $U_{\mathrm{MSD}}(S')\le U_{\mathrm{MSD}}(S)$.
\end{proof}

\begin{remark}
Note the asymmetry the two bound formulas encode: for SPSF/FPSF the ``deviation''
factor $|p_{\mathrm{base}}-P(h=1\mid X_S)|$ is bounded by a \emph{universal} constant
because a conditional probability always lies in $[0,1]$ regardless of $S$ — only the
\emph{weight} $P(X_S)$ needs to shrink monotonically. For MSD there is no such fixed
base rate to compare against; instead the two coverages $P_{D_1}(X_S),P_{D_2}(X_S)$
themselves box in every possible value of $\mathrm{SD}$ of any refinement, and the
bound is simply the larger of the two coverages, with no extra multiplicative factor.
Both are exactly the two cases described informally: ``$P(\text{subgroup})\cdot
\max(1-p_{\mathrm{base}},p_{\mathrm{base}})$'' and ``$\max$ of the subgroup's
probability under either distribution.''
\end{remark}

\section{Canonical generation order}

The search must also avoid re-deriving the same subgroup from several different
smaller conjunctions, and must know how to prune it in every such case at once. The
following resolves this at the level of \emph{how candidates are generated}, so that
no post-hoc bookkeeping is needed.

Fix an arbitrary but fixed order $a_1<a_2<\dots<a_d$ on $A$ — any order will do; it is
only ever used to break generation ties. For $S\neq\emptyset$ define
\[
\mathrm{top}(S):=\max\{i : a_i\in\mathrm{attr}(S)\}.
\]
In words, $\mathrm{top}(S)$ is simply \emph{the position, in the fixed order, of
whichever attribute already used by $S$ comes last} — informally, ``how far along the
attribute order $S$ has already reached.'' It is an index into $A$, not the size of
$S$: e.g.\ if $a_1<a_2<a_3<a_4$ and $S=\{a_1\!=\!v,\,a_3\!=\!w\}$, then $|S|=2$ but
$\mathrm{top}(S)=3$, because the attributes used by $S$ are $a_1$ and $a_3$, and $a_3$
is the later of the two. We set $\mathrm{top}(\emptyset):=0$ (the empty subgroup has
not reached any attribute yet, so the ``next'' attribute available to it is $a_1$).

Using this, define the \emph{canonical parent} of $S$ as the subgroup obtained by
dropping $S$'s \emph{last-indexed} literal:
\[
\pi(S):=S\setminus\{\ell\},\quad \ell \text{ the literal of $S$ on attribute } a_{\mathrm{top}(S)} .
\]
Continuing the example, $\pi(S)=\{a_1\!=\!v\}$ (we drop the literal on $a_3$, the
attribute achieving $\mathrm{top}(S)=3$, and keep the literal on $a_1$).

\begin{lemma}[Unique generation]\label{lem:tree}
The map $\pi$ makes $(S_A,\pi)$ a rooted tree with root $\emptyset$, whose children are
\[
\mathrm{children}(T)=\bigl\{\,T\cup\{a_j\!=\!v\}\ :\ a_j\in A,\ j>\mathrm{top}(T),\ v\in\mathrm{Val}(a_j)\,\bigr\}.
\]
In words: a subgroup's children are obtained by adding one literal on an attribute
that comes \emph{strictly after} $a_{\mathrm{top}(T)}$ in the fixed order — never on
an attribute earlier than or equal to $\mathrm{top}(T)$. In the running example,
$T=\{a_1\!=\!v,a_3\!=\!w\}$ has $\mathrm{top}(T)=3$, so its children add a literal on
$a_4$ (since $4>3$); they never add one on $a_2$, since $2<3=\mathrm{top}(T)$. This is
exactly what keeps the tree loop-free: $\{a_1\!=\!v,a_2\!=\!u,a_3\!=\!w\}$ is
\emph{not} a child of $T$ — it is reached only via its own canonical parent
$\{a_1\!=\!v,a_2\!=\!u\}$ (drop $a_3$, the last-indexed literal), a different node at
the same level as $T$. So the two size-3 refinements of $\{a_1{=}v\}$ that use $a_2$
or $a_3$ next are never confused with one another; each has exactly one route into the
tree. Consequently every $S\in S_A$ has a unique generating path from the root, i.e.\
is constructed exactly once, by exactly one parent.
\end{lemma}
\begin{proof}
Write the literals of any non-empty $S$ in increasing order of attribute index,
$a_{i_1}<\dots<a_{i_k}$. This exhibits exactly one chain
$\emptyset\prec\{a_{i_1}\!=\!v_1\}\prec\{a_{i_1}\!=\!v_1,a_{i_2}\!=\!v_2\}\prec\dots\prec S$
in which each step appends the literal whose attribute index is currently largest.
$\pi$ removes precisely the last element of \emph{this} chain, so $\pi$ and ``append
the literal with the next-largest index'' are mutually inverse operations; hence the
chain above is the \emph{only} root-to-$S$ path, and $\pi(S)$ is well defined and
unique.
\end{proof}

This is exactly what lets the algorithm ``neatly'' handle candidates that would
otherwise arise from several different conjunctions: with attributes always appended
in the same fixed global order, there simply is no second parent. If a subgroup $S$ is
pruned, every subgroup that $S$'s literal set is a subset of is either a strict
descendant of $S$ in this tree (and is therefore never even instantiated, because
generation only proceeds from survivors) or is not a superset reachable through $S$ at
all — but then it has its own, different, canonical parent $S''\ne S$ with
$S''\preceq$ it, which is checked against the very same bound independently and
correctly on its own path. No candidate is ever checked twice, and no candidate can
silently evade the check via ``another route,'' because it has exactly one route.

\section{Algorithm}

Recall $\mathrm{top}(S)$ is just the index (in the fixed attribute order
$a_1<\dots<a_d$) of the last attribute already used by $S$, with
$\mathrm{top}(\emptyset):=0$; ``$j>\mathrm{top}(S)$'' below means ``any attribute not
yet used by $S$ and coming after all attributes $S$ already uses.''

\begin{algorithm}[H]
\caption{Level-wise bound-and-prune search for $\argmax_{S\in S_A}V(S)$}
\label{alg:main}
\begin{algorithmic}[1]
\Require attributes $a_1<\dots<a_d$ with value sets $\mathrm{Val}(a_i)$; a violation
function $V$ and a valid anti-monotone bound $U$ for it (Propositions~\ref{prop:spsf}--\ref{prop:msd})
\Ensure $\tau=\max_{S\in S_A}V(S)$ and a subgroup $S^\star$ attaining it
\State $\tau \gets 0$; \quad $S^\star \gets \textsc{None}$
\State $\mathrm{Frontier} \gets \{\{a_j\!=\!v\}\ :\ a_j\in A,\ v\in\mathrm{Val}(a_j)\}$
       \Comment{all size-1 subgroups}
\While{$\mathrm{Frontier}\neq\emptyset$}
  \State $\mathrm{NextFrontier}\gets\emptyset$
  \ForAll{$S\in\mathrm{Frontier}$} \Comment{one full level: fixed $|S|=k$, single pairs before pairs before triples...}
    \If{$U(S)\le\tau$} \Comment{\emph{bound check, before computing $V(S)$}}
      \State \textbf{continue} \Comment{prune $S$: skip it, and generate none of its refinements}
    \EndIf
    \State $v \gets V(S)$ \Comment{only survivors pay for the exact violation}
    \If{$v>\tau$}
      \State $\tau\gets v$;\quad $S^\star\gets S$
    \EndIf
    \ForAll{$a_j\in A$ with $j>\mathrm{top}(S)$, \ $v'\in\mathrm{Val}(a_j)$} \Comment{canonical children only}
      \State $\mathrm{NextFrontier}\gets\mathrm{NextFrontier}\cup\{S\cup\{a_j\!=\!v'\}\}$
    \EndFor
  \EndFor
  \State $\mathrm{Frontier}\gets\mathrm{NextFrontier}$
\EndWhile
\State \Return $(\tau,S^\star)$
\end{algorithmic}
\end{algorithm}

\begin{center}
\resizebox{\textwidth}{!}{%
\begin{tabular}{lllll}
\hline
Measure & weight $w(S)$ used in $U$ & precomputed once & $U(S)$ & $V(S)$\\
\hline
SPSF & $P(X_S)$ & $p_{\mathrm{base}},\ \bar\delta_{\mathrm{SP}}$ &
$P(X_S)\bar\delta_{\mathrm{SP}}$ & $P(X_S)\,\bigl|p_{\mathrm{base}}-P(h=1\mid X_S)\bigr|$\\
FPSF & $P(X_S,y=0)$ & $p_{\mathrm{base},0},\ \bar\delta_{\mathrm{FP}}$ &
$P(X_S,y{=}0)\bar\delta_{\mathrm{FP}}$ & $P(X_S,y{=}0)\,\bigl|p_{\mathrm{base},0}-P(h=1\mid X_S,y{=}0)\bigr|$\\
MSD & --- & --- & $\max\bigl(P_{D_1}(X_S),P_{D_2}(X_S)\bigr)$ & $\bigl|P_{D_1}(X_S)-P_{D_2}(X_S)\bigr|$\\
\hline
\end{tabular}}
\end{center}

\section{Proof of validity}

\begin{theorem}[Validity]\label{thm:main}
Assume $S_A$ is finite (finitely many attributes, each with a finite domain). Let
$(\tau,S^\star)$ be the pair returned by Algorithm~\ref{alg:main}, instantiated with any
$(U,V)$ satisfying (B1)--(B2) — in particular any of Propositions
\ref{prop:spsf}--\ref{prop:msd}. Then
\[
\tau=\max_{S\in S_A}V(S) \qquad\text{and}\qquad V(S^\star)=\tau.
\]
That is, the algorithm never discards a subgroup that could have improved on the true
optimum: it is guaranteed to terminate having found a (globally) maximally violating
subgroup, together with the correct maximal value.
\end{theorem}

\begin{proof}
Because $S_A$ is finite, $\tau^\ast:=\max_{S\in S_A}V(S)$ exists; fix any
$S^{\mathrm{opt}}\in\argmax_{S\in S_A}V(S)$.

\smallskip
\emph{Invariant: $\tau\le\tau^\ast$ at every point of the execution.}
$\tau$ starts at $0\le\tau^\ast$ (as $V\ge0$) and is only ever reassigned (line~9) to a
value $v=V(S)$ for some subgroup $S$ that was actually evaluated; every such $v$
satisfies $v\le\tau^\ast$ by definition of $\tau^\ast$. So $\tau$ is only ever set to
values $\le\tau^\ast$, hence remains $\le\tau^\ast$ throughout, in particular at
termination.

\smallskip
It remains to show that the terminal $\tau$ satisfies $\tau\ge\tau^\ast$; together with
the invariant this forces $\tau=\tau^\ast$.

By Lemma~\ref{lem:tree}, $S^{\mathrm{opt}}$ has a unique generating chain
\[
\emptyset=S_0\prec S_1\prec\dots\prec S_m=S^{\mathrm{opt}}
\]
in which each $S_i$ ($i\ge1$) is placed on the Frontier at level $i$ precisely because
$S_{i-1}=\pi(S_i)$ survived the check at line~6 when it was processed (this is how
children are generated at lines 12--14). Two exhaustive cases:

\smallskip
\noindent\textbf{Case 1: every $S_i$, $1\le i\le m$, survives its check at line~6.}
Then $S_m=S^{\mathrm{opt}}$ itself is processed and not pruned, so line~8 computes
$v=V(S^{\mathrm{opt}})=\tau^\ast$ exactly. At that moment the current $\tau$ satisfies
$\tau\le\tau^\ast=v$ (invariant), so the condition on line~9 fires and $\tau$ is raised
to (at least) $\tau^\ast$. By the invariant this means $\tau=\tau^\ast$ from that point
on, hence also at termination.

\smallskip
\noindent\textbf{Case 2: some $S_i$ in the chain is pruned at line~6.}
Let $i^\ast\in\{1,\dots,m\}$ be the smallest such index, and let $\tau_{i^\ast}$ denote
the value of $\tau$ at the moment $S_{i^\ast}$ is checked; by definition of pruning,
\[
U(S_{i^\ast})\le\tau_{i^\ast}.
\]
Since $S_{i^\ast}\preceq S_m=S^{\mathrm{opt}}$, Lemma~\ref{lem:transfer} gives
\[
\tau^\ast=V(S^{\mathrm{opt}})\le U(S_{i^\ast})\le\tau_{i^\ast}.
\]
As $\tau$ is non-decreasing throughout execution (it is only ever replaced by a larger
value at line~9, otherwise left unchanged), $\tau_{i^\ast}\le\tau_{\mathrm{final}}$.
Combining, $\tau^\ast\le\tau_{\mathrm{final}}$.

\smallskip
In either case $\tau_{\mathrm{final}}\ge\tau^\ast$; with the invariant this gives
$\tau_{\mathrm{final}}=\tau^\ast$. Finally, $S^\star$ is only ever updated together
with $\tau$ (lines~8--9), so $V(S^\star)=\tau$ holds at every point of execution and
in particular at termination, giving $V(S^\star)=\tau_{\mathrm{final}}=\tau^\ast=
\max_{S\in S_A}V(S)$.
\end{proof}

\begin{remark}
Case~2 of the proof is the formal version of the claim in the problem statement: the
moment a subgroup is pruned because its bound no longer exceeds the running best,
\emph{the true global optimum has already been found} — it is sitting in $\tau$ at
that very instant (the chain of inequalities $\tau^\ast\le U(S_{i^\ast})\le
\tau_{i^\ast}\le\tau^\ast$ forces $\tau_{i^\ast}=\tau^\ast$). Pruning therefore never
loses the optimum; it only ever stops the search from re-deriving a value it has
already obtained by other means.
\end{remark}

\begin{corollary}[Certificate use]
If Algorithm~\ref{alg:main} terminates with $\tau\le\gamma$ for the SPSF (resp.\ FPSF)
instantiation, then $h$ satisfies $\gamma$-SPSF (resp.\ $\gamma$-FPSF) with respect to
$S_A$ — \emph{no} subgroup in $S_A$, evaluated exactly, needed to be checked to know
this.
\end{corollary}

\begin{corollary}[One algorithm, three measures]
By Corollary~1.1 of the unifying paper, running Algorithm~\ref{alg:main} with the pair
$(U_{\mathrm{MSD}},\mathrm{SD})$ on $(D_{h=1},D_{h=0})$ returns an $S^\star$ that also
maximizes $\mathrm{SPSF}(h,\cdot)$ over $S_A$ (the returned \emph{values} differ by the
$S$-independent factor $P(h=1)P(h=0)$, but the arg-max subgroup coincides); likewise,
by $\mathrm{FPSF}(h,S)=P(y=0)\cdot\mathrm{SPSF}_{D\mid y=0}(h,S)$, running it with
$(U_{\mathrm{FP}},\mathrm{FPSF})$ is the same search as running the SPSF instantiation
on the sub-population $\{y=0\}$. One implementation of Algorithm~\ref{alg:main},
parameterized by the $(U,V)$ pair in the table above, therefore suffices for all three
notions.
\end{corollary}

\section{Remarks}

\begin{itemize}
\item \textbf{Zero-mass subgroups.} If $P(X_S)=0$ (SPSF) or $P(X_S,y=0)=0$ (FPSF),
then $U(S)=0\le\tau$ automatically (as $\tau\ge0$ always), so such $S$ is pruned on
its very first check — consistent with the fact that a subgroup with no support can
have no support in any of its refinements either (Lemma~\ref{lem:cov}), so nothing of
interest is lost.
\item \textbf{Order within a level is irrelevant to correctness.} The proof of
Theorem~\ref{thm:main} places no requirement on the order in which subgroups of the
same level are processed, nor on breadth-first vs.\ depth-first traversal of the tree
of Lemma~\ref{lem:tree}; any order yields the same final $\tau$. Processing
high-violation subgroups earlier tends to raise $\tau$ sooner and thus prune more, but
this affects only running time, never correctness.
\item \textbf{Ties.} If several subgroups attain $\tau^\ast$, the argument in
Theorem~\ref{thm:main} applies to \emph{any} of them, so the algorithm is guaranteed to
return \emph{some} maximizer, though possibly not a specific one fixed in advance.
\item \textbf{Where the savings come from.} For SPSF/FPSF, $U(S)$ needs only the cheap
marginal (or joint-with-$y{=}0$) coverage of $S$, so a pruned subgroup never requires
computing the conditional probability $P(h=1\mid X_S)$ at all — exactly the ``before
computing the violation for each group'' savings requested. For MSD, computing
$U_{\mathrm{MSD}}(S)$ already requires both $P_{D_1}(X_S)$ and $P_{D_2}(X_S)$, which is
essentially all $\mathrm{SD}(S)$ needs too; there the main benefit of pruning is not
skipping per-node work but skipping entire pruned \emph{subtrees} of exponentially many
descendants, via Lemma~\ref{lem:transfer}.
\item \textbf{Worst-case complexity is unaffected.} Auditing for SPSF/FPSF over
conjunctions is NP-hard in the worst case (Theorem 3.3 of Kearns et al.\
\cite{kearns18}), and pruning cannot repeal that; the guarantee here is soundness
(the result is always exact when the algorithm finishes), not polynomial running time.
In practice, however, whole branches of the exponential lattice are cut as soon as
their bound drops below the incumbent, which is the entire empirical motivation for
using this scheme instead of naive enumeration.
\end{itemize}

\begin{thebibliography}{9}
\bibitem{msd} Maximum Subgroup Discrepancy (MSD): $\mathrm{MSD}(\mu,\nu;P)=\sup_{S\in
S_P}|\mu(\chi_S)-\nu(\chi_S)|$, together with its sample-complexity and MIO
formulation.
\bibitem{kearns18} Statistical Parity (SP) and False Positive (FP) Subgroup Fairness,
and the auditing-hardness result (Theorem 3.3).
\bibitem{unify} Unification of SPSF and MSD for detecting the most unfair subgroup
(Definition of Subgroup Discrepancy, Theorem 1, Corollary 1.1, and Eq.~(1) relating
FPSF to SPSF on the negative class).
\end{thebibliography}

\end{document}
